\section{Methods}
\label{sec:methods}

Every design choice in this chapter was forced by a specific failure of the obvious alternative, and
the chapter is organised as that chain. We need a quantity to measure (data and representation),
a ruler that cannot be gamed (metrics), manipulations that isolate drug use from everything that
imitates it (controls), and a way to ask where inside the model the failure sits (probing). Only
then do the target constructions make sense, because each is an attempted fix for a defect the
earlier sections have already localised. The two most consequential choices come last: how the
data are split, which decides what any number can claim, and what the model is compared against,
which decides what any number means. A reader who finishes the chapter should be able to predict
the shape of the results chapter from it.

\subsection{Data}
\label{sec:methods-data}

The corpus is Tahoe-100M~\cite{tahoe100m}: drug-treated single-cell transcriptomes across cancer
cell lines, with plate-matched DMSO controls. Metadata supplies cell line (Cellosaurus ID), drug,
dose (\texttt{drugname\_drugconc}), plate, and a fine-grained mechanism-of-action annotation
(\texttt{moa-fine}). All analysis is restricted to the $946$ genes in the intersection of the L1000
landmark set~\cite{l1000} with Tahoe's measured genes, and expression is library-normalised to
counts-per-10k and $\log1p$-transformed.
% AUDITOR NOTE: the production normalisation is log10(1 + x) after CP10K
% (tahoe_c2s_preprocess_endcell_v2.py:214), not natural log1p. One word wrong here; anyone
% re-implementing from this text gets a different transform.

One representational pitfall is worth recording because it is silent. Tahoe's per-cell
\texttt{genes} field holds \emph{vocabulary token identifiers}, not positional indices into the gene
metadata table. Mapping them positionally scrambles gene identity quietly --- the values exceed the
metadata row count, so the error surfaces only as an out-of-bounds index at the tail. All expression
construction here maps \texttt{gene\_metadata['token\_id']} $\rightarrow$ gene symbol $\rightarrow$
panel index, and the construction is validated by a PCA cell-line probe on the resulting panel,
which reaches $0.864$ against a chance rate of $0.02$.

\paragraph{Unit of analysis.} A \emph{condition} is a (drug, cell line, plate, dose) tuple, and the
physical treatment assignment is the \emph{sample} --- one treated well. That distinction, cosmetic
here, becomes load-bearing in Section~\ref{sec:methods-holdout}: one well holds roughly $49$ cell
lines dosed together, so a condition is an observation nested inside an assignment, not an
assignment. Plates are numbered within cell line rather than globally, so the batch unit is the
pair (cell line, plate) and never the plate label alone --- a distinction that matters for every
within-plate comparison below. Dose is carried as a molar concentration resolved from
\texttt{drugname\_drugconc}, not as a unitless float: an earlier pipeline kept only the numeric
part, which collides $1\,\mu$M with $1\,$nM, and a direct audit of all $289$ treatment samples found
$289$ usable molar doses with zero unit collisions --- every drug in this atlas states its doses in
a single unit, so the defect was a demonstrated possibility that did not materialise. Combination
samples are dropped rather than analysed as their first component.

\paragraph{Scale.} Table~\ref{tab:scale} states the quantities every $n$ in this thesis should be
located against. Two populations are used and they are not interchangeable. The broad
characterisation of the task (Section~\ref{sec:res-blind}) streams a wide sample of the atlas; the
residual arm (Sections~\ref{sec:res-encoding} onward) works from a smaller cached subset built once
and reused, so that every arm sees identical cells.

\begin{table}[h]
\centering
\begin{tabular}{lr}
\toprule
\multicolumn{2}{l}{\textbf{Task characterisation} (streamed sample)} \\
\midrule
conditions characterised & 6,628 \\
\quad of which statistically active against DMSO & 3,647 / 5,014 tested (72.7\%) \\
\quad identifiable in principle (ceiling $\geq 0.8$) & 3,064 (46.2\%) \\
\quad with a plate-mate closer than their own replicate & 78.7\% \\
median cells per condition & 44 \\
median \#DEGs per condition (BH $q<0.05$) & 0 \\
median SNR (effect $\div$ replicate noise) & 0.75 \\
drugs with a real dose series & 218 / 351 (62.1\%) \\
drugs seen in $\geq 3$ cell lines & 227 \\
\midrule
\multicolumn{2}{l}{\textbf{Residual-frame cache} (rebuilt targets, Section~\ref{sec:methods-residual})} \\
\midrule
cells & 620,564 (600,000 treated, 20,564 control) \\
(cell line, plate) groups & 198 \\
drugs / cell lines & 106 / 50 \\
conditions with $\geq40$ treated and $\geq20$ control cells & 6,617 \\
\quad by split after inventory filters (train / held-out wells / held-out drugs) & 4,999 / 900 / 711 \\
\quad retained by the data-resolved reliability line & 2,523 of 4,999 train (50.5\%) \\
training examples / well-held-out validation examples & 147,710 / 3,600 \\
\midrule
\multicolumn{2}{l}{\textbf{Evaluation sets}} \\
\midrule
tier 1 (seen conditions) / tier 2 (unseen drugs) / tier 3 (unseen combinations) & 4,325 / 2,188 / 99 \\
conditions scored per tier in the generation evaluations & 400 \\
tier-2 drugs in the expression-frame $\NIR$ benchmark (within plate) & 606 \\
conditions scored in the residual-frame evaluation & 570 (41 cell lines, 90 drugs) \\
\bottomrule
\end{tabular}
\caption{The populations behind every $n$ quoted in this thesis. Sources:
\texttt{drug\_biology\_atlas.json}, \texttt{eval\_endcell\_cpu.json},
\texttt{nir\_sameplate.json}, the rebuilt target report and the residual-frame evaluation. Note
that the 72.7\% active figure is over the 5,014 conditions with enough cells for the permutation
test, not over all 6,628.}
\label{tab:scale}
\end{table}

Three evaluation tiers, defined by what the training set contained, are used for the broad
characterisation. \emph{Tier 1} holds conditions whose (drug, cell line) pairing was seen in
training, so it measures fit rather than generalisation. \emph{Tier 2} holds conditions whose
\emph{drug} never appeared, which is Leave-Drugs-Out in the vocabulary of
Section~\ref{sec:lit-dreval}. \emph{Tier 3} holds pairings whose drug and cell line were each seen
but never together --- Leave-Pairs-Out. Tier 3 is small ($99$ conditions) and is reported only
where it is adequately powered; the residual arm builds its own, larger holdout, and the unit that
holdout is keyed at turns out to be the most consequential choice in this chapter
(Section~\ref{sec:methods-holdout}).

\subsection{Model and representation}
\label{sec:methods-model}

A cell is written as a \emph{cell sentence}: its expressed panel genes in descending order of
expression, terminated by \ENDCELL. The prompt supplies cell line, drug, dose, mechanism, and the
untreated control cell's own sentence; the response is the treated cell's sentence. Two tokens,
\ENDCELL{} and --- for the residual arm --- \DOWN{}, are registered as atomic special tokens; left
unregistered, \DOWN{} splits into sub-words and the up/down boundary ceases to be a clean symbol.

The backbone is \texttt{C2S-Scale-Pythia-1b-pt}~\cite{c2sscale}, a C2S-Scale checkpoint built on
the Pythia-1b decoder~\cite{pythia}, cold-started for every arm from the same base checkpoint with
identical hyperparameters (1 epoch, learning rate $10^{-5}$, gradient accumulation 16, maximum
length 8192, bf16), so that in each experiment the training target or objective is the only
variable.

What that checkpoint has already seen bears directly on how the held-out splits should be read. It
is pretrained on over 50 million human and mouse cells drawn from 825 single-cell datasets in
CELLxGENE and the Human Cell Atlas --- an \emph{atlas} corpus, carrying no chemical or genetic
perturbation data, and the released model card lists cell and tissue prediction, conditioned
generation and gene-set conversion among its capabilities but not perturbation prediction. The base
model therefore reaches our fine-tuning having never seen a drug-treated cell, which is what allows
the unseen-drug split to mean what it claims: a drug held out here is unseen by the whole pipeline,
not merely by its final stage.

\subsection{Metrics}
\label{sec:methods-metrics}

The first thing this thesis needs is a ruler that cannot be gamed, because the field's default one
can be --- and Section~\ref{sec:res-metrics} exhibits the exploit rather than asserting it. The
default is a correlation between predicted and true \emph{change from control}, restricted to the
genes that actually move. Writing $x$ for the control profile, $y$ for the treated truth and
$\hat{y}$ for the prediction, let
\begin{equation}
  S_K \;=\; \operatorname*{arg\,top-}K_{\,g} \; \lvert y_g - x_g \rvert
\end{equation}
be the $K$ genes with the largest true change. Then
\begin{equation}
  \DEdr \;=\; \operatorname{corr}\!\big( \left(\hat{y} - x\right)\big|_{S_K}, \;
                                          \left(y - x\right)\big|_{S_K} \big).
\end{equation}
We report $K=50$ with Pearson correlation throughout; $K \in \{20, 100, 200\}$ and the Spearman
variant were computed and move the numbers by less than $0.02$ except at $K=200$, where the
top-$K$ set starts to include genes that do not move.

Three properties of this definition, each innocent in isolation, compose into the failure. First,
$S_K$ is selected using the \emph{truth}, so every predictor is scored on the genes the condition
actually moved --- the exam questions are chosen from the answer key. Second, treated cells in this
atlas regress toward a shared response, so a prediction that simply reverts toward the population
centre correlates with $y - x$ \emph{by construction}: no drug knowledge is required to score well,
only knowledge of where the centre is. Third, the fix is available in the same breath ---
\textbf{partial-}$\DEdr$ regresses the control profile out of both arguments before correlating,
which removes the reversion route and is therefore the honest version of the same quantity.
Section~\ref{sec:res-metrics} shows what each version scores a predictor that contains no
information at all.

A cell sentence lists only expressed genes, so a predicted profile leaves some panel genes
unranked, and any rank-based metric needs a convention for the absent tail. Two were used
throughout and reported side by side: \texttt{worst} assigns every absent gene the last rank $P$,
and \texttt{francesca} assigns them a fixed middle rank $P/2$. The choice is immaterial --- spike-in
$\DEdr$ reads $0.952$ against $0.954$ under the two, and no baseline or ceiling in this thesis
moves by more than $0.01$ --- so all figures are quoted under \texttt{worst} and the convention is
not mentioned again. For a subset-free alternative we also compute $\paneltau$, Kendall's $\tau$
between predicted and true rank over the full $946$-gene panel, which is invariant to the
absent-gene convention by construction.

The metric this thesis relies on is the Normalised Inverse Rank, introduced by Wu et al.\ in
PerturBench~\cite{perturbench} (there called the ``transposed-rank'' metric) and identified by
Miller et al.~\cite{miller} as one of the few consistently calibrated metrics across 14 Perturb-seq
datasets. For a prediction $\hat{y}$ of condition $i$, with truths $\{y_j\}$ over the comparison
set,
\begin{equation}
  \NIR(\hat{y}, i) \;=\; \frac{1}{|J|}\sum_{j \in J}
    \Big[\mathbf{1}\{\sigma(\hat{y}, y_i) > \sigma(\hat{y}, y_j)\}
       + \tfrac{1}{2}\,\mathbf{1}\{\sigma(\hat{y}, y_i) = \sigma(\hat{y}, y_j)\}\Big],
\end{equation}
the fraction of \emph{other} conditions whose truth the prediction resembles less than its own,
with ties counted at one half. Two properties do work that no other candidate does. Chance is
$0.50$ \emph{by construction}, not empirically: ``this model is at chance'' is a statement, not a
comparison against a fitted null, and that is what makes the negative results in this thesis
quotable without a control arm. And because the generic response is present in every condition's
truth in equal measure, it cancels in the comparison --- $\NIR$ is a \emph{discrimination} metric
and only the drug-specific difference decides it, which is precisely the component the reversion
artifact exploits. The tie term is not cosmetic: in residual space
(Section~\ref{sec:methods-residual}) a predictor of ``the average drug response'' is exactly the
zero vector, every similarity against it is zero, and a strict inequality scores it $0.000$ rather
than the $0.500$ it deserves.

\paragraph{Which $\NIR$ variant, and why it must be stated.} Our instruments compute $\NIR$ under
two similarity functions and store both: \texttt{nir\_expr}, Euclidean distance in expression space,
and \texttt{nir\_rank}, the same construction over rank vectors. \textbf{Every $\NIR$ figure in this
thesis is \texttt{nir\_expr}}, and the choice is load-bearing rather than incidental: under
\texttt{nir\_rank} the tier-2 aggregate places \emph{every} arm below chance, including the noise
ceiling itself ($0.380$, against the model's $0.313$ and a drug-agnostic linear map's $0.285$). A
metric on which a real replicate of a condition fails to be retrieved above chance is not measuring
what its name suggests, and it is not the metric the calibration analysis of
Section~\ref{sec:res-metrics} endorsed --- the $\DRF$ result is computed on the expression-space
variant. We report the calibrated one and name it here so the stored artifacts cannot be misread.

% fig-nir.tex -- NIR as a discrimination score. Goes with the metric definition.
% No data dependency.
\begin{figure}[htbp]
\centering
\begin{tikzpicture}[
  x=1cm, y=1cm, font=\small,
  lbl/.style={font=\scriptsize, text=black!70},
  strong/.style={font=\scriptsize, text=black!90},
  own/.style={circle, fill=black!85, inner sep=2pt},
  other/.style={circle, draw=black!55, fill=white, inner sep=1.9pt},
  ar/.style={-{Stealth[length=4.5pt]}, black!65},
]

\node[lbl, anchor=south] at (2.05,0.44) {other drugs in the same comparison set};
\node[strong, anchor=south] at (6.30,0.44) {own truth};

\fill[black!10, rounded corners=1pt] (0.55,0.00) rectangle (6.30,0.34);
\foreach \x in {0.95, 1.75, 2.60, 3.55, 4.45, 5.45} {\node[other] at (\x,0.17) {};}
\node[own] at (6.30,0.17) {};
\node[other] at (7.15,0.17) {};
\node[other] at (7.85,0.17) {};
\node[anchor=west, font=\small] at (8.35,0.17) {$\NIR = \tfrac{6}{8}$};

\draw[ar] (0.45,-0.26) -- (8.15,-0.26);
\node[lbl, anchor=north] at (4.30,-0.38)
  {similarity of the prediction to each condition's truth $\longrightarrow$};

\draw[decorate, decoration={brace, mirror, amplitude=3.5pt}, black!55]
  (0.55,-0.80) -- (6.30,-0.80);
\node[strong, anchor=north] at (3.40,-0.96) {the shaded fraction \emph{is} the score};

\node[lbl, anchor=west] at (0.45,-1.50) {own truth ranked first $\Rightarrow \NIR \to 1$};
\node[lbl, anchor=west] at (0.45,-1.84) {own truth ranked at random $\Rightarrow \NIR = 0.5$};
\node[strong, anchor=west] at (0.45,-2.18)
  {every similarity \emph{identical} $\Rightarrow \NIR = 0.5$ by the half-tie rule, not $0$};

\end{tikzpicture}
\caption[NIR is a discrimination score]{\textbf{$\NIR$ asks whether a prediction resembles its own
condition's truth more than it resembles the truths of the other drugs in the same comparison set.}
The generic response --- the stress and cell-cycle program every compound triggers --- is present in
every truth in equal measure, so it cancels in the comparison and only the drug-specific difference
decides. Two consequences follow, and both matter downstream. Chance at $0.50$ is a statement about
the \emph{comparison set}, not about the model: conditions are grouped by (cell line, plate) in the
expression frame and by cell line in the residual frame, and the two are not interchangeable. And a
predictor whose similarities are all identical --- the zero vector, which is what ``predict the
average drug response'' becomes in residual space --- must score $0.50$ rather than $0$, which is why
ties are counted at one half rather than dropped.}
\label{fig:nir}
\end{figure}


Finally, the metrics themselves need grading, and for that we port the framework of Miller
et al.~\cite{miller}, asking how well each \emph{metric} separates a positive from a negative
control:
\begin{equation}
  \DRF(m) \;=\; \frac{m(\textrm{pos}) - m(\textrm{neg})}{m(\textrm{perfect}) - m(\textrm{neg})},
\end{equation}
with $\textrm{pos}$ an interpolated-duplicate noise ceiling --- their positive control, which
interpolates a technical duplicate toward the mean baseline with per-gene weights given by
differential-expression significance --- and $\textrm{neg}$ a drug-agnostic leave-one-out mean. A
negative \DRF{} means the metric scores the uninformative baseline \emph{above} the noise ceiling:
it is inverted, and rewards the generic gene ordering rather than the drug. Three hardening choices
matter for how the result can be read. Tied expression values receive \emph{average} ranks, because
breaking ties by array index makes a tied gene's rank depend on its position in the panel. The
\DRF{} ratio is recomputed \emph{whole} inside every bootstrap resample, because its denominator is
estimated from the same data as its numerator and holding it fixed understates the interval. And
the family of five metrics is controlled simultaneously with Holm's procedure, because the claim
``only $\NIR$ is calibrated'' is a claim about all five at once.

Adopting this framework is deliberate. Miller et al.\ use it to argue that reports of deep models
failing to beat uninformative baselines reflect miscalibrated evaluation rather than model failure
(Section~\ref{sec:lit-dispute}). Evaluating our model on the metric \emph{their} analysis endorses
means the result reported in Section~\ref{sec:res-blind} cannot be dismissed on the grounds their
own paper raises.

\subsection{Controls}
\label{sec:methods-controls}

The negative results in this thesis are only as strong as the controls that make them
interpretable. Each of the following exists because a specific objection would otherwise be fatal,
and each is stated with the failure it closes.

Drug and plate are partially confounded by Tahoe's design: each drug is assayed on its own
plate(s), with roughly 4--12 drugs per plate and plate-matched controls. Discrimination grouped by
cell line alone therefore lets the \emph{batch} signature stand in for drug identity. Measured
directly, by scoring the same tier-2 conditions with and without the plate held constant:

\begin{table}[h]
\centering
\begin{tabular}{lrr}
\toprule
Predictor (drug-agnostic) & cross-plate & within plate \\
\midrule
mean baseline & 0.399 & \textbf{0.161} \\
control-copy & 0.551 & 0.510 \\
\midrule
noise ceiling & 0.625 & 0.575 \\
fraction of drugs at chance & 43\% & 49\% \\
\bottomrule
\end{tabular}
\caption{Expression-frame $\NIR$. A drug-agnostic mean baseline more than doubles its score when the
plate is free to vary, which is what plate leakage looks like. The ceiling falls once the plate is
held constant because the batch signature was inflating it too.}
\label{tab:plateleak}
\end{table}

The mean baseline is the clearest demonstration: it contains no drug information whatsoever, and it
scores $0.399$ when the comparison set may span plates against $0.161$ when it may not. Every
discrimination instrument therefore gained a \texttt{--same\_plate\_only} mode and all discrimination
results in this thesis were re-run within plate. Conclusions are unchanged; some magnitudes shrink.

A separate and stronger observation, which belongs to the drug-blindness argument rather than to this
one, is that \emph{within} plate and on the identifiable stratum a control-copy baseline reaches
$0.766$ against the model's $0.768$ (Section~\ref{sec:res-blind}). That is not a statement about
plate leakage --- it is a statement about the model.

\paragraph{The scramble ablation, and what its comparator must satisfy.} The central manipulation
leaves the prompt byte-identical except that the drug name and mechanism are replaced by another
drug's, while the control cell and the scoring truth are unchanged. Because both arms share the
same control cell, control- and batch-derived leakage cancels in the \emph{difference}
$\gap = \NIR_{\textrm{model}} - \NIR_{\textrm{scramble}}$. This is the only manipulation that
isolates drug use, and every drug-use claim in this thesis is anchored to it rather than to a
comparison against a drug-agnostic predictor.

Two constraints on the swap partner are required for a gap measured this way to mean what it says,
and both exist because the unconstrained pool fails in a specific way. The partner must be a
\emph{different drug}: allowing the same drug at another dose or plate substitutes the drug for
itself, an unchanged output is then correct rather than blind, and the contamination concentrates
in the most similar stratum because a drug's own other doses are usually its nearest neighbours.
And the partner must be on the \emph{same plate}: residuals retain plate structure, so a
cross-plate partner differs from the target in batch as well as in drug, and the gap absorbs both.

A scramble that swaps in a drug with a \emph{similar} true response is also a weak test, and this is
a real risk here because same-mechanism drugs are barely more similar than random pairs (within-MoA
distance $2.322$ versus between-MoA $2.377$, ratio $0.977$), so ``different mechanism'' does not
imply ``different response''. Two remedies are used: a \emph{swap-distance sweep} that turns the
single test into a curve of $\gap$ against the response distance $\lVert y_A - y_B \rVert$ between
the real and swapped-in drug; and a \emph{stratified scramble} with three defined partners per
condition chosen by residual cosine within the cell line --- \texttt{near} (most similar),
\texttt{orth} ($\cos\approx0$) and \texttt{opposite} (most anti-correlated). The strata make an
assumption visible that a single comparator hides: the gap has two parts, how much naming the truth
helps and how much naming a lie hurts, and only the first is wanted. A usable null is therefore a
stratum whose own score sits at chance --- which, Section~\ref{sec:res-generalisation} shows, is
true of exactly one of the three, and not the one an evaluation would naturally default to.

% fig-scramble.tex -- the scramble ablation and why its difference is leak-immune.
% No data dependency.
% Layout note: the prompt boxes are ~2.2cm tall once the shaded spans are set, which is far more
% than a four-line block looks like it should need. Every band below them is placed against that
% measured height, not against an estimate.
\begin{figure}[htbp]
\centering
\begin{tikzpicture}[
  x=1cm, y=1cm, font=\small,
  lbl/.style={font=\scriptsize, text=black!70},
  strong/.style={font=\scriptsize, text=black!90},
  promptbox/.style={draw=black!35, rounded corners=2pt, inner sep=5pt, fill=black!2},
  ar/.style={-{Stealth[length=4.5pt]}, black!65},
]
\newcommand{\spanhl}[1]{\colorbox{black!14}{#1}}

\node[lbl, anchor=west] at (0.15,0.00) {real prompt};
\node[promptbox, anchor=north west] (p1) at (0.15,-0.26) {%
  \renewcommand{\arraystretch}{0.9}%
  \begin{tabular}{@{}l@{~~}l@{}}
  \ttfamily\scriptsize Cell line:    & \ttfamily\scriptsize MCF7 \\
  \ttfamily\scriptsize Drug:         & \ttfamily\scriptsize \spanhl{Lapatinib} \\
  \ttfamily\scriptsize Mechanism:    & \ttfamily\scriptsize \spanhl{EGFR inhibitor} \\
  \ttfamily\scriptsize Control cell: & \ttfamily\scriptsize ACTB GAPDH \ldots \\
  \end{tabular}};

\node[lbl, anchor=west] at (0.15,-2.86) {scrambled prompt --- only the two shaded spans differ};
\node[promptbox, anchor=north west] (p2) at (0.15,-3.12) {%
  \renewcommand{\arraystretch}{0.9}%
  \begin{tabular}{@{}l@{~~}l@{}}
  \ttfamily\scriptsize Cell line:    & \ttfamily\scriptsize MCF7 \\
  \ttfamily\scriptsize Drug:         & \ttfamily\scriptsize \spanhl{Vorinostat} \\
  \ttfamily\scriptsize Mechanism:    & \ttfamily\scriptsize \spanhl{HDAC inhibitor} \\
  \ttfamily\scriptsize Control cell: & \ttfamily\scriptsize ACTB GAPDH \ldots \\
  \end{tabular}};

\node[draw=black!35, rounded corners=2pt, inner sep=6pt, align=center, fill=black!3]
  (truth) at (9.30,-2.55) {truth for\\\textbf{Lapatinib}\\{\scriptsize drawn once}};

\draw[ar] (p1.east) .. controls (6.6,-1.35) and (7.6,-1.85) .. (truth.north west);
\draw[ar] (p2.east) .. controls (6.6,-4.20) and (7.6,-3.30) .. (truth.south west);

\node[anchor=north, font=\small] at (9.30,-3.70)
  {$\gap = \NIR_{\text{real}} - \NIR_{\text{scrambled}}$};

\draw[ar, dashed, black!45] (1.20,-6.10) -- (8.10,-6.10);
\node[strong, anchor=south, align=center] at (4.65,-6.02)
  {the shared control cell enters \emph{both} arms, so control- and batch-derived};
\node[strong, anchor=north, align=center] at (4.65,-6.20)
  {leakage cancels in the difference};

\end{tikzpicture}
\caption[The scramble ablation]{\textbf{The scramble replaces the drug name and its mechanism and
nothing else.} Both arms are conditioned on the same control cell and scored against the same truth,
so any leakage through the control cell or the plate enters both arms in equal measure and cancels in
the difference $\gap$. This is why $\gap$, rather than a comparison against a drug-agnostic
predictor, is the quantity every drug-use claim in this thesis is anchored to: the drug-agnostic
comparison is not leak-immune, and a drug-agnostic mean baseline scores $\NIR = 0.399$ when the
comparison set may span plates against $0.161$ when it may not, precisely because it inherits that
leakage. A refinement matters here and is
developed in Section~\ref{sec:methods-controls}: if the swapped-in drug happens to have a
\emph{similar} true response then an unchanged output is correct rather than blind, so the swap
partner is chosen from three defined similarity strata and the evidence is the \emph{gradient} across
them rather than any single value.}
\label{fig:scramble}
\end{figure}


The scramble replaces the drug name \emph{and} its mechanism together, so a non-null gap
establishes sensitivity to the combined drug-and-mechanism prompt without attributing it to either
span. A \emph{field decomposition} costs one extra pair of generation arms and swaps exactly one
span of the instruction line to the opposite-signature partner, leaving every other byte identical:
the drug name with the mechanism kept, the mechanism with the name kept, and both together. One
guard is essential rather than tidy. If the swap partner happens to share the original's mechanism,
then swapping only the mechanism produces a prompt \emph{identical} to the model's own, and the arm
scores a gap of exactly zero \emph{by construction}. Read naively that is indistinguishable from
``the model ignores the mechanism'', which is the very claim under test. Such conditions are
dropped rather than scored, which is why the mechanism arm carries fewer conditions than the others.

\paragraph{Two further drug-isolating instruments.} The drug-blindness claim rests on four
instruments rather than one, because they fail in different ways and agreement across them is worth
more than any single number. Two are defined above ($\NIR$ and the scramble); the other two are
forced-choice grading and output invariance.

In \emph{forced-choice grading}, a prediction is scored against its own condition's truth and
against another drug's truth from the same comparison set, and the instrument records how often the
own truth wins. Chance is $0.50$ by construction. The ceiling is the same forced choice with a real
held-out replicate substituted for the prediction, which bounds what any predictor could achieve at
this noise level; it lands between $0.67$ and $0.83$ depending on cells per condition.

\emph{Output invariance} asks whether the drug token changes the generated text at all, without
reference to any truth. For a fixed control cell, the model generates twice from the \emph{same}
prompt and once from each of several prompts differing only in the drug, and the instrument compares
\begin{equation}
  \textrm{gap} \;=\; \underbrace{\overline{\textrm{sim}}(\textrm{same prompt, resampled})}_{\textrm{sampling noise alone}}
  \;-\; \underbrace{\overline{\textrm{sim}}(\textrm{prompts differing in the drug})}_{\textrm{noise plus any drug effect}},
\end{equation}
with similarity measured as top-$100$ Kendall $\tau$ and as Jaccard overlap. A positive gap means two
generations for the same drug resemble each other more than generations for different drugs do,
i.e.\ the drug token moves the output. A gap of zero means swapping the drug perturbs the output no
more than resampling at the same temperature does. Reported over $8$ contexts $\times$ $12$ drugs at
temperature $0.8$. This instrument is the weakest of the four --- it can only detect whether the
output moves, not whether it moves \emph{correctly} --- but it is also the one least able to be
confounded, since it needs no ground truth at all.

\paragraph{Channel gate.} For an unseen drug, prediction is only possible through a property the drug
\emph{shares} with drugs seen in training, so the reachable channels can be enumerated and each one
measured in the same frame. For a channel $X$, a drug's residual is predicted as the mean residual of
the \emph{other} drugs in the same cell line that share property $X$ with it --- the construction of
the mechanism baseline, generalised. Because it uses only drugs that were in training, it is a fair
contest in every split, unlike a same-drug lookup, which on an unseen drug is an oracle.

The gate's estimand requires the channel arm and its null to differ in exactly one respect, and two
disciplines enforce it. Every channel is paired with a \textbf{plate-matched random null}: the same
\emph{number} of partner drugs, drawn at random from the same cell line, \emph{and} the same split
of partners across the target's own plate and other plates, because mechanism- and target-related
drugs are co-plated more often than chance --- a null matched on count alone would differ from the
channel arm in plate membership as well as in the channel, which is two differences where the
estimand allows one. A \textbf{co-plating diagnostic is printed before any verdict}, so the size of
the correction is a number rather than an argument, and where a plate-matched null spans too few
cell lines the channel is reported \textsc{untestable} rather than closed --- ``we could not build
the control'' must never be written as ``closed''. Coverage is likewise reported before any
verdict, because a channel annotated on few drugs cannot be closed by a null result: it was never
tested.

The bound this establishes is on \emph{similarity transfer} through each channel, not on any model.
A model could in principle combine channels or learn a non-smooth structure-to-response map that a
nearest-neighbour average cannot express. A leave-one-drug-out regression over all channels jointly
addresses the first. The second is bounded by sample size rather than by argument: at a few hundred
drugs, learning a non-smooth map is not feasible, so at this $n$ the lookup is effectively the bound.

\paragraph{Is the interaction learnable?} The transfer coefficient says how \emph{large} the
drug$\times$cell-line interaction is; it says nothing about whether anything could learn it. If a cell
line modulated drug responses in a consistent direction, that direction would be estimable from the
line's own cells and a conditional model would have a concrete target. Writing
$\kappa(d,c) = \resid(d,c) - \hat{\beta}(d)$, the test compares $\cos(\kappa(d,c), \kappa(d',c))$ for
different drugs in the \emph{same} cell line against the same quantity across \emph{different} cell
lines.

$\hat{\beta}$ must be estimated \textbf{leave-one-cell-line-out}. Estimated over all lines it would
contain $\resid(d,c)$ itself, and subtracting a mean containing its own term induces a $-1/(m-1)$
correlation that biases the signal arm downward \emph{by construction} --- which would make a
structured interaction look idiosyncratic, the exact error the test exists to avoid. Both arms are
disattenuated by $\kappa$'s own split-half reliability, since $\kappa$ is a difference of two noisy
quantities and is therefore noisier than the residual it comes from.

\paragraph{Uncertainty.} Conditions are not independent, and they are not independent in two
crossed ways at once. Two conditions from the same treated well share a physical assignment and a
batch --- one well holds about $49$ cell lines dosed together --- and two conditions from the same
cell line share the biology. Treating conditions as independent gives intervals that are too
narrow; here, about $1.3$ to $1.5$ times too narrow. Neither grouping alone captures the
dependence, and the tempting shortcut is actively misleading: clustering on the assignment unit
alone uses the \emph{larger} cluster count ($289$ wells against $\approx 50$ lines) and would
\emph{tighten} every interval while sounding more rigorous. All intervals in this thesis are
therefore two-way cluster-robust (Cameron--Gelbach--Miller),
\begin{equation}
  V \;=\; V_{\textrm{line}} \;+\; V_{\textrm{well}} \;-\; V_{\textrm{intersection}},
\end{equation}
where the intersection of the two groupings is the single condition, so $V_{\textrm{intersection}}$
is the ordinary independent-sampling variance; when the estimator goes non-positive in a finite
sample the wider one-way variance is used and the branch recorded. Because a cluster-robust
variance is asymptotic in the number of \emph{clusters}, not observations, the reference
distribution is Student's $t$ with degrees of freedom one less than the smaller cluster count ---
which matters where an arm spans only $25$ wells, and moves a borderline verdict. For statistics
computed over \emph{pairs} of conditions, such as the transfer coefficient, neither grouping
partitions the observations --- a pair belongs to two nodes at once --- so those intervals use the
Fafchamps--Gubert dyadic estimator, under which two pairs are dependent whenever they share either
member. This is the structure DrEval~\cite{dreval} identifies in this field more broadly, citing
Hurlbert's pseudoreplication~\cite{hurlbert}: hundreds of thousands of cells resolving to a few
hundred conditions over $50$ cell lines.

\paragraph{Generation validity.} Every generation-based evaluation reports the \ENDCELL{}/\DOWN{}
completion rate, the fraction of emitted symbols that are valid panel genes, and the duplicate rate,
before any score is quoted. This is not hygiene: a $600$-token generation cap silently truncated
$26\%$ of generations in an early run and \emph{halved} the measured effect. Every contrast is
additionally generated with a fixed per-(condition, arm) seed, so the two arms of a paired
comparison are not drawn from different points of one global stream.

\subsection{Mechanistic probing}
\label{sec:methods-probe}

To distinguish ``cannot represent the drug'' from ``represents but does not use it'', four
instruments are needed, and again each exists because its predecessor fails in an instructive way.

Per-layer \emph{linear probes} on the residual stream at the last prompt token, predicting drug
identity, establish whether the information is present at all. But the raw drug subspace is partly
a batch subspace --- probing dose and plate from \emph{inside} it recovers dose at $0.90$--$0.95$ ---
so a context subspace is built from the union of cell line, plate and dose, and the drug axis is
orthogonalised against it; every causal claim uses the purified subspace. Before any ``ablating the
drug does nothing'' result is trusted, a \emph{removal gate} requires the subspace to contain the
drug: the fraction of above-chance probe accuracy removed by projecting it out must exceed $0.8$,
without which the null is unfalsifiable. The \emph{causal test} then reads log-probabilities of the
teacher-forced true response clean, then re-reads them with a forward hook projecting the subspace
out at every position, taking the mean KL divergence over response tokens; measuring in logit space
rather than by generating removes the sampling-noise floor that made an earlier steering experiment
inconclusive.

\paragraph{The identity test, and two designs that failed before it.} Ablation shows a subspace
carries functional variance; it cannot show that the subspace's \emph{drug-specific content} is read.
Only a substitution can. Two attempts at one were retracted, and since the failures were instructive
they are stated rather than hidden.

The first overwrote drug A's code with drug B's inside a purified subspace and compared the resulting
KL against a random vector of identical norm drawn in the \emph{full} residual stream --- so almost
all of the control's energy lay in directions the intervention never touched. The second drew the
random vector inside the subspace, which appeared to fix it and did not: the hook ablates before
adding, so the delivered perturbation is (added $-$ removed), and because the injected class mean was
uncentred it shared its global-mean component with the thing removed. The substitution arm was then
a near-restore and the control arm a genuine displacement. Simulated with no model in the loop, the
control delivers $5.3\times$ the perturbation in $100\%$ of draws. Both retracted results pointed the
same way, and that way was geometry rather than biology.

The design used here removes the ablation from the identity test. It injects a \emph{displacement}
$P(\mu_B - \mu_A)$ --- in which the global mean cancels exactly --- and scores a \emph{contrast},
\begin{equation}
  \log P\!\left(r_B \mid \textrm{prompt}_A\right) \;-\; \log P\!\left(r_A \mid \textrm{prompt}_A\right),
\end{equation}
so that any shift common to both responses, including the entire normaliser, subtracts out. What
survives is the model's preference between drug B's own held-out response and drug A's. The headline
control is a second real drug displacement $P(\mu_C - \mu_A)$, rescaled to identical norm: same
anchor, same construction, wrong drug. An isotropic in-subspace direction is retained only as an
off-manifold damage diagnostic, never as the comparison, because a random direction disrupts the
model in ways an on-manifold class-mean displacement does not.

Subspaces and class means are fit on one split, $r_B$ drawn from a second disjoint split, and all
scoring done on a third. Uncertainty is a cluster bootstrap over (A,B) pairs; equivalence requires
the interval to lie inside a pre-registered margin \emph{and} contain zero; and a layer is excluded
where the cell-line positive control has no dynamic range, since a saturated instrument cannot
support either verdict.

\subsection{Target constructions}
\label{sec:methods-targets}

The three target families below are attempts to fix the target \emph{content} while leaving its
encoding alone; Section~\ref{sec:res-epsilon} shows why that family of fix could not work, and the
ladder construction is what makes that demonstration conclusive rather than anecdotal.

\emph{Consensus.} Cells grouped by (drug, cell line, dose) and averaged into one pseudobulk
sentence, with every original prompt retained so only the response changes.

\emph{Optimal transport.} For each (cell line, plate) group, a coupling $\pi$ between control
and treated cells minimising $\sum_{ij}\pi_{ij}\lVert x_i - y_j\rVert^2$ under uniform marginals,
solved by Sinkhorn iteration~\cite{cuturi} on $K = \exp(-C/\varepsilon)$. Couplings are computed in
the released Tahoe scVI latent~\cite{scvi} ($n_{\textrm{latent}} = 10$), joined per cell by barcode
rather than re-encoded --- a re-encode produced a degraded latent (cell-line probe $0.11$ versus
$0.975$ for the shipped latent). Targets are built \emph{decoder-free}, as convex combinations of
real treated cells in expression space, so they always lie on the data manifold.

The regularisation parameter defines a ladder whose endpoints are experiments we had already run:
$\varepsilon\rightarrow\infty$ gives a uniform coupling and hence the treated mean (consensus);
$\varepsilon\rightarrow 0$ gives near-hard assignment, approximately a single state-matched treated
cell. This is what makes the sweep conclusive rather than anecdotal.

\subsection{Residual targets}
\label{sec:methods-residual}

The residual target asks the model to predict only what makes \emph{this} drug different. For each
condition (drug, cell line, plate, dose) with at least $40$ treated and $20$ plate-matched control
cells:
\begin{equation}
  \resid(d,c) \;=\; \underbrace{\big(\bar{y}_{\textrm{treated}} - \bar{y}_{\textrm{control}}\big)}_{\shift(d,c)}
  \;-\; \underbrace{\frac{1}{|D_c|}\sum_{d' \in D_c} \shift(d',c)}_{\textrm{generic}(c)} .
\end{equation}
The control subtraction removes the cell's own state; the mean-over-drugs subtraction removes the
\emph{generic drug program} --- the stress and cell-cycle response common to all compounds, which
Section~\ref{sec:res-metrics} shows accounts for essentially all of the model's apparent skill. What
remains is what makes \emph{this} drug different. Four properties of how the generic is estimated
decide whether that remainder can be trusted, and each was forced by a measured failure.

First, \textbf{the split precedes the fit}. An earlier build estimated the generic and the
reliability filter over \emph{every} condition and assigned the holdout afterwards, so each held-out
target was defined partly by other held-out conditions. The pipeline is therefore staged ---
inventory (metadata only), split (metadata only, before any pseudobulk exists), fit (the generic is
handed the training keys explicitly), transform (held-out conditions are projected through the
fitted generic and never enter it) --- and a poison test stands behind the staging: corrupting every
held-out expression vector must leave the training targets byte-identical, and corrupting a training
condition must change them.

Second, \textbf{the scope is the plate, not the cell line}. The original argument for cell-line
scope was reproducibility --- $62\%$ of conditions reproduce across a half-split at cell-line scope
against $19\%$ at plate scope, on the reasoning that a plate holds too few drugs to estimate a mean
and subtracting a noisy one injects noise. That comparison does not hold up: both figures were
computed with a \emph{shared} control pseudobulk subtracted from both halves, which makes the halves
correlate through their common control error, and the inflation is not neutral between scopes ---
it favours cell-line scope precisely because the cell-line generic is itself less noisy. Measured
with disjoint control halves, the two are comparable: $20\%$ at plate scope against $16\%$ at
cell-line scope. Two further diagnostics settle it in plate's favour and are reported in
Section~\ref{sec:res-transfer}: at cell-line scope the dose-transfer ordering inverts and the
shared-minus-split control bias reads $+0.254$, against $-0.000$ at plate scope. A plate group with
fewer than three other training drugs yields no generic, and the condition is dropped and counted
rather than silently falling back to the contaminated frame.

Third, \textbf{the mean is over drugs, not conditions}, and \textbf{leaves the condition's own drug
out}. A drug measured at five doses would otherwise contribute five times the weight of a
single-dose drug. And once the generic is train-only, an asymmetry appears that had not existed
before: a training condition's generic would contain its own drug while a held-out condition's would
not, so the two splits would be scored against differently defined targets and the generalisation
gap would partly measure that difference. The generic for every condition is therefore ``the mean
over training drugs other than this one, within scope'' --- identical in definition on both sides.

Fourth, \textbf{the reliability line is resolved from the data, not inherited}. A condition is kept
for training only if its residual reproduces across a half-split, $\cos(\resid_A, \resid_B)$ above
a threshold --- with each half measured against its \emph{own} control half, so the two share
nothing. The threshold itself was the single largest determinant of the training set, and it had
been inherited ($0.2$) from an era when reliability was measured with shared controls --- a
different quantity on a different scale, under which the mean is $\approx 0.6$ rather than
$+0.127$. Rather than pick a new number, the builder reports the whole retention curve against a
null that pairs half A of one condition with half B of a \emph{different} condition --- same
encoding, same dimensionality, same noise, no shared biology --- and the threshold is set at the
null's $95$th percentile: a condition is kept when its two halves agree more than two unrelated
conditions' halves do. At the measured null ($95$th percentile $+0.109$) this retains $2{,}523$ of
$4{,}999$ training conditions ($50.5\%$). The held-out sets are \emph{not} filtered by this
criterion: dropping held-out conditions for failing their own reproducibility bar would select the
evaluation set on its own outcome, so noise in retained training conditions costs learning
efficiency, not validity.

What this fraction means is stated carefully, because it is easy to over-read. The half-split
agreement is a \emph{within-well sampling-precision} criterion: it asks whether two disjoint cell
subsets of the \emph{same} well give the same residual. It is not biological reproducibility ---
$286$ of $287$ distinct (drug, dose) treatments in this atlas occur in exactly one well, so no
independent-well replicate exists to measure that with --- and the $50.5\%$ is a statement about
how much drug-specific signal survives pseudobulk noise at this cell budget, not about how much
biology is real.

\paragraph{Encoding and scoring.} The target string is
\texttt{<up genes> \DOWN{} <down genes> \ENDCELL}, $100$ genes per block, ordered by signed residual
magnitude; sign is biology, so it is given a symbol rather than being folded into a magnitude
ranking. Predicted and true residuals are both mapped to a signed rank vector (top-$k$ up positive,
top-$k$ down negative, weight $1/\log_2(\textrm{rank}+1)$) so a generated sentence and a continuous
truth are compared like for like; similarity is cosine and the score is $\NIR$ over other drugs in
the same cell line.

% fig-frame.tex -- the residual frame and the encoding fork.
% No data dependency. Annotations from RESULTS/target_divergence.json, cellline scope:
%   full 34.59 / shift 111.42 / residual 118.25 ; ceilings 0.7420 / 0.7420 / 0.7432.
% Layout notes:
%   * two subtraction ROWS, not an L-shape (the L put the generic term on a band that collided
%     with its own caption).
%   * the caption for the generic sits BELOW both rows, not between them -- between them it ran
%     under the divider and into the right-hand column.
%   * row labels are centred on profiles spaced 1.65cm apart; at 1.45cm "generic(c)" and
%     "residual r(d,c)" touch.
\begin{figure}[htbp]
\centering
\begin{tikzpicture}[
  x=1cm, y=1cm, font=\small,
  bar/.style={fill=black!55, draw=none},
  faintbar/.style={fill=black!22, draw=none},
  gene/.style={draw=black!35, rounded corners=1pt, inner sep=1.8pt, font=\ttfamily\scriptsize},
  shared/.style={gene, fill=black!5, text=black!45},
  differs/.style={gene, fill=black!16, text=black, draw=black!70, thick},
  lbl/.style={font=\scriptsize, text=black!70},
  strong/.style={font=\scriptsize, text=black!90},
  op/.style={font=\normalsize, text=black!75},
]
\newcommand{\prof}[4]{%
  \foreach \h [count=\i from 0] in {#3} {
    \fill[#4] ($(#1,#2)+(\i*3.0pt,0)$) rectangle ++(2.3pt,\h pt);
  }%
}

% ================================================ LEFT: two subtractions
\node[lbl]    at (0.52,3.34) {treated};
\node[lbl]    at (2.17,3.34) {control};
\node[strong] at (3.90,3.34) {shift $\shift(d,c)$};
\prof{0.10}{2.62}{17,12,9,14,7,10,5,8}{bar}
\node[op] at (1.48,2.92) {$-$};
\prof{1.75}{2.62}{10,8,6,11,4,7,3,6}{bar}
\node[op] at (3.13,2.92) {$=$};
\prof{3.40}{2.62}{7,4,3,3,3,3,2,2}{bar}

\node[strong] at (0.52,1.94) {shift};
\node[lbl]    at (2.17,1.94) {generic$(c)$};
\node[strong] at (3.95,1.94) {residual $\resid(d,c)$};
\prof{0.10}{1.22}{7,4,3,3,3,3,2,2}{bar}
\node[op] at (1.48,1.52) {$-$};
\prof{1.75}{1.22}{6,4,3,2,3,2,2,2}{faintbar}
\node[op] at (3.13,1.52) {$=$};
\prof{3.40}{1.22}{1,0,0,1,0,1,0,0}{bar}

\node[lbl, anchor=north, align=center] at (2.30,0.82)
  {generic$(c)$ is the mean over the \emph{other}\\drugs in this cell line};

\draw[black!25] (5.05,0.10) -- (5.05,3.60);

% ================================================ RIGHT: the encoding fork
\node[strong, anchor=west] at (5.35,3.44) {the \emph{same} measurement, written two ways};

\node[lbl, anchor=west] at (5.35,2.96) {\textbf{(a)} genes ranked by \emph{expression}};
\node[anchor=west] at (5.35,2.52) {%
  \tikz[baseline]{\node[shared]{ACTB};}\;\tikz[baseline]{\node[shared]{GAPDH};}\;%
  \tikz[baseline]{\node[differs]{CDKN1A};}\;\tikz[baseline]{\node[shared]{TPT1};}\;$\cdots$};
\node[anchor=west] at (5.35,2.06) {%
  \tikz[baseline]{\node[shared]{ACTB};}\;\tikz[baseline]{\node[shared]{GAPDH};}\;%
  \tikz[baseline]{\node[differs]{HMOX1};}\;\tikz[baseline]{\node[shared]{TPT1};}\;$\cdots$};
\node[strong, anchor=west] at (5.35,1.58)
  {only \textbf{34.6 / 200} tokens differ \;(\textbf{17\%})};

\node[lbl, anchor=west] at (5.35,1.00) {\textbf{(b)} genes ranked by \emph{signed residual}};
\node[anchor=west] at (5.35,0.56) {%
  \tikz[baseline]{\node[differs]{CDKN1A};}\;\tikz[baseline]{\node[differs]{MDM2};}\;%
  \tikz[baseline]{\node[gene, fill=black!3]{[DOWN]};}\;\tikz[baseline]{\node[differs]{CCNB1};}\;%
  $\cdots$};
\node[strong, anchor=west] at (5.35,0.08)
  {\textbf{118.2 / 200} differ \;(\textbf{59\%})};

% ================================================ the invariant
\draw[decorate, decoration={brace, amplitude=4pt}, black!55] (0.10,-0.55) -- (11.20,-0.55);
\node[strong, anchor=north] at (5.65,-0.73)
  {and the replicate-retrieval ceiling is \textbf{unchanged}: $0.742 \rightarrow 0.743$};

\end{tikzpicture}
\caption[The residual frame and the encoding fork]{\textbf{The information was always present; the
tokenisation discarded it.} Left: the drug-specific residual is built by two subtractions. The
control removes the cell's own state, leaving the \emph{shift}; the mean over the other drugs in the
same cell line removes the \emph{generic} program that every compound triggers, leaving what makes
this drug different. Right: that same measurement, written two ways. Ranking genes by expression (a)
puts housekeeping genes first, so only $34.6$ of the top $200$ tokens differ between two drugs in the
same context and cross-entropy has almost nothing to learn drug identity from; ranking by signed
residual (b) raises that to $118.2$. The bracket is the point of the figure --- the
replicate-retrieval ceiling is essentially identical under both encodings, so no information has been
added. What changed is how much of it the token sequence exposes to the loss.}
\label{fig:frame}
\end{figure}


\subsection{Generalisation design}
\label{sec:methods-holdout}

This is the choice on which every generalisation claim in the thesis stands or falls, and it is
forced by the physical design of the atlas rather than by statistics. Tahoe applies one drug at one
dose to one well containing $\approx 49$ cell lines simultaneously --- measured: $289$ treatment
samples, a mean of $48.6$ cell lines per well (median $49$, range $43$--$50$), and $286$ of $287$
distinct (drug, dose) treatments occurring in exactly one well. Three consequences follow, and they
apply to every analysis of this atlas, not only this one.

A drug cannot be present in cell line A and absent from cell line B at the same dose, because the
two lines are physically the same well. So a $(\textrm{drug}, \textrm{cell line})$ holdout --- the
field's standard Leave-Pairs-Out split, and the split an earlier version of this work used ---
leaves the held-out condition's treated well sitting in the training set through its other
$\approx 48$ cell lines. Measured against the holdout this thesis previously used, $124$ of the
$268$ treated wells in this cache placed a held-out condition and a training condition in the same
physical sample. What that estimand measures is \emph{within-atlas matrix completion}, not
independent generalisation, and no processing of the targets changes it, because it is a property
of what was randomised.

Holding out \emph{whole wells} is the only assignment-respecting choice the design offers, and it
answers a different and narrower question: generalisation to an \emph{unseen treatment well of a
seen drug} --- a new physical preparation, not a new cellular context. Fixed-dose cross-context
transfer, present in one line and absent in another, is structurally unidentifiable in this atlas;
this is a property of Tahoe's design, not of our pipeline, and no rebuild fixes it. The split used
throughout the residual arm is therefore keyed at the well:
\begin{description}
  \item[\texttt{train}] every condition of the remaining wells ($4{,}999$ conditions before the
        reliability line, $2{,}523$ after it).
  \item[\texttt{unseen\_combo} --- unseen wells of seen drugs] $36$ treated wells held out whole
        ($904$ conditions, $900$ after inventory filters). The holdout never takes a drug's last
        training well --- without that guard a two-well drug could silently become an unseen
        \emph{drug} while sitting in the unseen-combination arm, making the two arms measure the
        same thing. Zero wells contribute conditions to both sides of this split.
  \item[\texttt{unseen\_drug} --- Leave-Drugs-Out] every condition of $10$ held-out drugs ($714$
        conditions, $711$ after inventory filters). The model never sees the token. This is the
        \emph{control}: naming a drug the model knows nothing about versus naming something else can
        make no difference, so the measured effect there must be zero --- a non-zero value is an
        instrument fault, not a biological finding.
\end{description}

We do not build a Leave-Cell-Lines-Out (LCO) split here; Section~\ref{sec:limitations} argues it is
the natural next axis and explains why its interest is contingent on the transfer coefficient.

\paragraph{A warning we inherit with the vocabulary.} DrEval reports, citing~\cite{partin}, that
apparently good performance on pair-holdout splits can be reached simply by predicting each drug's
average response across the training cell lines. Our held-out-well estimand sits exactly in the
regime where that warning bites: the drug is seen, and its average is computable. This is the
reason the baseline ladder below is not optional: the drug-average predictor is not a formality to
be dismissed, it is the hypothesis under test.

\paragraph{Stratified sampling.} Evaluation samples conditions under an explicit per-split quota. A
uniform draw reproduces the pool's proportions, which starves precisely the split the experiment
exists to test: in a first run, $250$ constructed held-out combinations yielded only $33$ scored
conditions, purely as an artifact of the sampler. The residual-frame evaluation therefore scores
$200$ train, $250$ \texttt{unseen\_combo} and $120$ \texttt{unseen\_drug} conditions --- the latter
two being quotas of the $900$ and $711$ available, so the held-out verdicts in
Section~\ref{sec:res-generalisation} are powered on a subset of the conditions the split affords,
and their detection limits should be read with that in mind.

\subsection{Baselines}
\label{sec:methods-baselines}

Each baseline is a predictor of the drug-specific residual, scored through an identical pipeline ---
same truths, same signed-rank encoding, same comparison set.

\begin{table}[h]
\centering
\begin{tabular}{lll}
\toprule
Baseline & Information used & Role \\
\midrule
\texttt{drug\_lookup}   & this drug's mean residual, all other cell lines & the bar to beat \\
\texttt{drug\_lookup\_1} & this drug in \emph{one} other cell line & removes the averaging advantage \\
\texttt{moa\_lookup}    & other drugs sharing this drug's MoA, same cell line & mechanism-leak diagnostic \\
\texttt{control\_copy}  & none (predict the control cell) & must be $\approx 0.50$ \\
\texttt{generic}        & none (predict the average response) & must be $\approx 0.50$ \\
\texttt{random}         & none & floor \\
\bottomrule
\end{tabular}
\caption{The baseline ladder. \texttt{drug\_lookup} is the pre-registered win condition: a lookup
can reproduce a memorised average, so only a model that tailors the response to the control cell's
state can beat it.}
\label{tab:baselines}
\end{table}

Two disciplines make the ladder fair. The lookups are fitted from \textbf{training conditions
only}: fitted from the full cache, a lookup scoring a held-out condition reads residuals the model
was never shown --- an oracle, not a baseline --- and the oracle variants are reported alongside,
explicitly labelled, as a bound on how much the comparison was ever inflated rather than hidden.
And \texttt{drug\_lookup\_1} selects a cell \emph{line} and averages that line's conditions rather
than selecting a single condition, because $62\%$ of drugs are multi-dose and a single condition
would confound ``different cell line'' with ``different dose''. \texttt{moa\_lookup} is not an
oracle in any split, as it uses only drugs that were in training.

\subsection{Variance decomposition}
\label{sec:methods-vardecomp}

To separate the per-drug main effect from the drug$\times$cell-line interaction we model
\begin{equation}
  \resid(d,c) \;=\; \beta(d) \;+\; \kappa(d,c) \;+\; \epsilon,
  \qquad
  \Tcoef \;=\; \frac{\operatorname{var}\beta}{\operatorname{var}\beta + \operatorname{var}\kappa},
\end{equation}
where $\beta$ is constant across cell lines and $\kappa$ is the interaction --- the only component a
conditional model can win on.

\paragraph{This is DrEval's naive predictor, in transcriptomic space.} DrEval's
\texttt{NaiveMeanEffectsPredictor}~\cite{dreval} predicts
$\hat{y}_{ij} = \mu + (\mu^c_i - \mu) + (\mu^d_j - \mu) = \mu^c_i + \mu^d_j - \mu$: an additive
model of cell-line and drug main effects, with no interaction term. Our
$\textrm{generic}(c) + \beta(d)$ is precisely that quantity, lifted from a scalar potency to a
$946$-dimensional profile. It follows that $\kappa$ \emph{is} the term their naive model omits, and
that $1 - \Tcoef$ measures exactly the headroom available above it. DrEval finds that deep models
barely beat this additive baseline on drug sensitivity; the transfer coefficient asks whether, in
the transcriptomic setting, there is anything above it to win in the first place.

$\Tcoef$ is estimated by disattenuated cross-line correlation,
\begin{equation}
  \hat{\Tcoef} \;=\; \frac{\cos\!\big(\resid(d,c_1),\, \resid(d,c_2)\big)}
                           {\sqrt{\rho(d,c_1)\,\rho(d,c_2)}},
  \qquad
  \rho \;=\; \frac{2\cos(\resid_A, \resid_B)}{1 + \cos(\resid_A, \resid_B)},
\end{equation}
with $\rho$ the Spearman--Brown reliability of the full-sample residual. Dividing out the
reliability is essential --- measurement noise alone depresses the raw cross-line cosine and would
otherwise be read as interaction --- and the reliability entering it is a \emph{within-well}
split-half precision, which is optimistic relative to true biological reproducibility. That
optimism sits in the denominator, so the interaction share $1 - \Tcoef$ is, if anything,
understated.

Five controls accompany the estimate, each capable of changing the conclusion:
\begin{enumerate}
  \item \textbf{Control half-split.} Subtracting one shared control pseudobulk from both halves
        leaves control-estimation error common to $\resid_A$ and $\resid_B$, inflating the
        reliability and biasing $\Tcoef$ \emph{downward} --- manufacturing apparent interaction.
        Both builds are run and the bias reported ($+0.254$ at cell-line scope, $-0.000$ at plate
        scope).
  \item \textbf{Negative control, structure-matched.} A different drug in a \emph{different} cell
        line, breaking only the drug match while preserving the cross-line structure of the
        estimand; it must read $\approx 0$. A same-plate/different-drug control was tried first and
        is \emph{not} adequate: it differs from the estimand in two ways at once, and plate
        structure surviving a cell-line-scoped generic inflates it to $+0.478$ while leaving the
        estimand untouched --- which produced a false-alarm rejection before the matched control was
        built.
  \item \textbf{Simulated null.} At $\kappa = 0$ the estimator does not return $1.0$. A simulation
        calibrated to the measured noise level fixes the reference. In validation the raw cross-line
        cosine at $\kappa = 0$ is $0.734$, which read naively reports $27\%$ interaction where the
        truth is zero.
  \item \textbf{Dose strata.} Same drug, same cell line, different dose is not cross-line transfer;
        it is reported separately as a reference scale, with dose resolved to a \emph{molar}
        quantity --- an earlier cache carried the well identifier in the dose field, so what was
        labelled dose transfer compared physical samples rather than concentrations.
  \item \textbf{Leave-one-out generic.} The generic includes drug $d$, leaving a $-1/m$ self-term,
        unless the condition's own drug is excluded --- which, per
        Section~\ref{sec:methods-residual}, it now is.
\end{enumerate}
Validation against planted ground truth recovers $\Tcoef = 1.00$, $0.70$ and $0.40$ to within
$0.005$, and the different-drug negative control returns $-0.001$.

Finally, an exact energy decomposition of the shift into residual, generic and their cross term ---
the three shares sum to one by construction --- reports what fraction of the total perturbation
response the drug-specific frame covers: the number that scopes every claim in this thesis, and one
that has not previously been computed for this task.
